Large Language Model (LLM) evaluation results vary across reproductions, but variance sources interact, and their relative importance shifts with the benchmark format. No unified factorial framework exists to decompose these sources simultaneously across benchmarks. We apply Generalizability theory (G-theory) in a crossed design with seven facets across four models and four benchmarks, totaling 1.27 million evaluation records. The variance structure is format-conditional: on multiple-choice benchmarks, item selection accounts for 42--49\% of total variance, while on the free-form benchmark GSM8K, model and prompt template each contribute 12--15\% with no single dominant source. Model$\times$item interaction (33.6\%) is the hidden second-largest component, larger than all methodology main effects combined. Under binary scoring, BF16 and FP16 numerical precisions are functionally interchangeable ($<$0.1\% interaction). Multi-facet designs require 8.8$\times$ fewer items than single-condition designs to reach a generalizability coefficient of 0.80 (72 versus 631 items). We release an open-source G-theory pipeline and Decision study (D-study) budget calculator for evaluation design optimization.
